
\subsection{退化条件：\textup{(D.3)} 的证明}
\label{subsec:6-4-2}
\liftIndex{I029}

本小节利用第~\ref{subsec:3-8-9} 小节构造的算子
$\mathfrak p$\liftIndex{I388} 证明 \textup{(D.3)}。
回忆该小节定义了算子列
\[
\begin{aligned}
\mathfrak p_{\ell,k}\colon{}&
E_\ell C(M;\Lambda_{0,\mathrm{nov}})[2]
\otimes
B_k^{\mathrm{cyc}}C_{1+}(L;\Lambda_{0,\mathrm{nov}})[1]\\
&\longrightarrow S^*(M;\Lambda_{0,\mathrm{nov}}).
\end{aligned}
\]
参见命题~\ref{prop:3-8-78}。
对 $(\boldsymbol b,b)\in\mathcal M_{\mathrm{def}}(L)$，定义
\begin{equation}
\begin{aligned}
\mathfrak p_1^{\boldsymbol b,b}(x)
&=\mathfrak p_{0,1}^{\boldsymbol b,b}(x)\\
&:=\mathfrak p(e^{\boldsymbol b},e^bxe^b)
=\mathfrak p(e^{\boldsymbol b}\otimes e^bxe^b)
\end{aligned}
\tag{6.4.1}\label{eq:6-4-1}
\end{equation}
其中 $x\in C_2(L;\Lambda_{0,\mathrm{nov}})[1]$。
注意
\[
\mathfrak p_1^{\boldsymbol b,b}(x)
=\mathfrak p^{\boldsymbol b}(e^bxe^b),
\]
这里的算子 $\mathfrak p^{\boldsymbol b}$
已在第~\ref{subsec:3-8-9} 小节定义。
还回忆
\[
\delta_{(\boldsymbol b,b)}(x)
=\mathfrak q(e^{\boldsymbol b},e^bxe^b)
=\mathfrak q(e^{\boldsymbol b}\otimes e^bxe^b)
\]
（见定义~\ref{def:3-8-58}）。

\begin{lemma}{6.4.2}\label{lem:6-4-2}
设 $\delta_M$ 是命题~\ref{prop:3-8-78}
中 $S^*(M;\Lambda_{0,\mathrm{nov}})$ 上的上边界算子。
则有
\begin{align}
\mathfrak p_1^{\boldsymbol b,b}
&\equiv i_!\pmod{\Lambda_{0,\mathrm{nov}}^+},
\tag{6.4.3.1}\label{eq:6-4-3-1}\\
\mathfrak p_1^{\boldsymbol b,b}
\circ\delta_{(\boldsymbol b,b)}
+\delta_M\circ\mathfrak p_1^{\boldsymbol b,b}
&=0,
\tag{6.4.3.2}\label{eq:6-4-3-2}
\end{align}
其中 $i_!:H^*(L;\mathbb Q)\to
H^{*+n}(M;\mathbb Q)$ 是 Gysin 同态。
\end{lemma}

\begin{proof}
由于 $\boldsymbol b\equiv b\equiv0
\pmod{\Lambda_{0,\mathrm{nov}}^+}$，
有 $\mathfrak p_1^{\boldsymbol b,b}
\equiv\mathfrak p_1
\pmod{\Lambda_{0,\mathrm{nov}}^+}$。
因此 (6.4.3.1) 由 (3.8.87.1) 得出。
又因为
\[
\begin{aligned}
\widehat d^{\boldsymbol b}(e^bxe^b)
&=e^b\mathfrak m^{\boldsymbol b}(e^bxe^b)e^b\\
&=e^b\mathfrak q(e^{\boldsymbol b},e^bxe^b)e^b,
\end{aligned}
\]
由 (3.8.87.2) 可得
\[
\begin{aligned}
0
&=\mathfrak p^{\boldsymbol b}
\circ\widehat d^{\boldsymbol b,\mathrm{cyc}}(e^bxe^b)
+\delta_M\circ\mathfrak p^{\boldsymbol b}(e^bxe^b)\\
&=\mathfrak p^{\boldsymbol b}
\bigl(e^b\mathfrak q(e^{\boldsymbol b},e^bxe^b)e^b\bigr)
+\delta_M\circ\mathfrak p_1^{\boldsymbol b,b}(x).
\end{aligned}
\]
这就证明了 (6.4.3.2)。
\end{proof}

当 $L$ 经 bulk 形变后弱无障碍时，
使用算子
\[
\mathfrak p_1^{+\boldsymbol b,b}\colon
C_{1+}(L;\Lambda_{0,\mathrm{nov}})^+[1]
\longrightarrow S^*(M;\Lambda_{0,\mathrm{nov}}),
\]
它定义为
\begin{equation}
\mathfrak p_1^{+\boldsymbol b,b}(x)
=\mathfrak p_{0,1}^{+\boldsymbol b,b}(x)
:=\mathfrak p^+(e^{\boldsymbol b},e^bxe^b)
\tag{6.4.4}\label{eq:6-4-4}
\end{equation}
其中 $(\boldsymbol b,b)
\in\mathcal M_{\mathrm{weak,def}}(L)$。
这里的 $\mathfrak p^+$ 定义于
命题~\ref{prop:3-8-78}。
与引理~\ref{lem:6-4-2} 同理，
可证下述引理；其证明从略。

\begin{lemma}{6.4.5}\label{lem:6-4-5}
在上述情形下，
\begin{align}
\mathfrak p_1^{+\boldsymbol b,b}
&\equiv i_!\pmod{\Lambda_{0,\mathrm{nov}}^+},
\tag{6.4.6.1}\label{eq:6-4-6-1}\\
\mathfrak p_1^{+\boldsymbol b,b}
\circ\delta_{(\boldsymbol b,b)}^+
+\delta_M\circ\mathfrak p_1^{+\boldsymbol b,b}
&=0.
\tag{6.4.6.2}\label{eq:6-4-6-2}
\end{align}
这里 $\delta_{(\boldsymbol b,b)}^+(x)
=\mathfrak q^+(e^{\boldsymbol b},e^bxe^b)$，
参见定义~\ref{def:3-8-58}。
\end{lemma}

\noindent\textbf{定理 D 的 \textup{(D.3)} 之证明。}
置 $C=C(L;\Lambda_{0,\mathrm{nov}})$、
$\delta=\delta_{(\boldsymbol b,b)}$。
取 $x\in Z_r^{p,q}(\widehat C)$。
也就是说，设 $x=y+z$，其中
$\delta y\in\mathfrak F^{q+r-1}
\widehat C^{p+1}$，
$z\in\mathfrak F^{q+1}\widehat C^p$。
注意
\begin{equation}
\begin{aligned}
\delta_r([x])
&=[\delta y]
\in E_r^{p+1,q+r-1}(\widehat C)
=\frac{Z_r^{p+1,q+r-1}(\widehat C)}
{B_r^{p+1,q+r-1}(\widehat C)}\\
&=\frac{Z_r^{p+1,q+r-1}(\widehat C)}
{\bigl(
\delta(\mathfrak F^{q+1}\widehat C^p)
\cap\mathfrak F^{q+r-1}\widehat C^{p+1}
\bigr)+\mathfrak F^{q+r}\widehat C^{p+1}}.
\end{aligned}
\tag{6.4.7}\label{eq:6-4-7}
\end{equation}
由 (6.4.3.2) 有
\begin{equation}
\mathfrak p_1^{\boldsymbol b,b}(\delta y)
=-\delta_M\bigl(\mathfrak p_1^{\boldsymbol b,b}(y)\bigr).
\tag{6.4.8}\label{eq:6-4-8}
\end{equation}
由 (6.4.3.1)，
可以在 (6.3.18) 中选取 $\lambda_0$，使
\begin{equation}
\mathfrak p_1^{\boldsymbol b,b}-i_!
\equiv0\pmod{T^{\lambda_0}}.
\tag{6.4.9}\label{eq:6-4-9}
\end{equation}
因为
$\delta y\in\mathfrak F^{q+r-1}\widehat C^{p+1}
=T^{(q+r-1)\lambda_0}\widehat C^{p+1}$，
由 (6.4.8)、(6.4.9) 得
\begin{equation}
i_!(\delta y)
+\delta_M\bigl(\mathfrak p_1^{\boldsymbol b,b}(y)\bigr)
\in T^{(q+r)\lambda_0}
S^{p+n+1}(M;\Lambda_{0,\mathrm{nov}}).
\tag{6.4.10}\label{eq:6-4-10}
\end{equation}
因此可写 $\delta y=w_1+w_2$，
其中
$w_2\in\mathfrak F^{q+r}\widehat C^{p+1}$，
且
\begin{equation}
i_!(w_1)\in\operatorname{Im}\delta_M.
\tag{6.4.11}\label{eq:6-4-11}
\end{equation}
这里用到 $i_!$ 与 $\delta_M$
均不含 $T$ 的事实。\translatornote{原书实体页 386 在式 (6.4.10) 前
误引 (6.4.7)、(6.4.8)，
推得此式实际还须用紧前的 (6.4.9)；
式末又写 $C^{p+1}(M)$，
但 $i_!$ 升高次数 $n$，且靶为 $S^*(M)$。
此处依型别写作 $S^{p+n+1}(M)$，
见疑误 A-0374--A-0375。}

注意，借助庞加莱对偶，
Gysin 同态
$i_!:H^k(L;\mathbb Q)
\to H^{n+k}(M;\mathbb Q)$
与
$i_*:H_{n-k}(L;\mathbb Q)
\to H_{n-k}(M;\mathbb Q)$
相对应。
因此由 (6.4.7)、(6.4.11)
及 $K_r$ 的定义，
可得 $\delta_r([y])\in K_r$。
\translatornote{原书实体页 386 此处写 $K_{r+1}$；
但 (6.4.7) 的像属于 $E_r$，
而定理 D 将 $K_{r+1}$ 定义为 $E_{r+1}$ 的子群。
所证 (D.3) 也正是
$\operatorname{Im}\delta_r\subset K_r$，
见疑误 A-0373。}
性质 \textup{(D.3)} 得证。
\hfill$\square$

(6.2.15) 的证明与此类似。

\begin{remark}{6.4.12}\label{rem:6-4-12}
一般而言，当 $b_1\ne b_0$ 时，
\textup{(D.3)} 不适用于弗洛尔上同调
\[
HF\bigl((L;(\boldsymbol b,b_1)),
(L;(\boldsymbol b,b_0));
\Lambda_{0,\mathrm{nov}}\bigr).
\]
例如，令
$L=S^1=S^1\times\{0\}
\subset S^1\times S^1=M$。
若 $\boldsymbol b=0$、$b_1=0$、
$b_0=T\operatorname{PD}([pt])$，则
\[
\begin{aligned}
\delta_{b_1,b_0}
\bigl(\operatorname{PD}([S^1])\bigr)
&=\mathfrak m_2
\bigl(\operatorname{PD}([S^1]),b_0\bigr)\\
&=T\cdot\operatorname{PD}([pt]),
\end{aligned}
\]
而该元素不属于
\[
\ker\bigl(
i_!:H^1(L;\Lambda_{0,\mathrm{nov}})
\longrightarrow H^2(M;\Lambda_{0,\mathrm{nov}})
\bigr).
\]
\end{remark}
